\begin{proposition}[Quadratic sharpness of the pair caps]
\label{prop:pair-cap-extremal}
There are infinitely many triples $(p,H,\mathcal Z)$, where $p$ is an
odd prime, $p>H^2$, and
$\mathcal Z\subseteq\{0,\ldots,p-1\}$, with the following properties:
\begin{enumerate}
 \item $z\in\mathcal Z$ if and only if $p-1-z\in\mathcal Z$;
 \item $0,p-1\notin\mathcal Z$, and no two elements of $\mathcal Z$ are
 consecutive;
 \item for every $1\le h<p$,
 \[
  A_{\mathcal Z}(h)
  :=\#\{z:z,z+h\in\mathcal Z\}\le3(h-1);
 \]
 \item the ordered set $\mathcal Z$ has at least
 \[
  \frac{H^2}{7200}-4
 \]
 consecutive four-point windows of span at most~$H$, none of which
 contains a pair whose sum is $p-1$.
\end{enumerate}
Consequently, reflection, absence of adjacent points, and the individual
Ap\'ery pair caps alone cannot give a subquadratic bound for the number of
short off-center four-point windows.
\end{proposition}

\begin{proof}
Put $D=1$.
Let $m\ge1100$ be a multiple of~$100$, and put
\[
 R=\frac m{100},\qquad H=12m,qquad N=mR.
\]
For $0\le i<N$, define
\[
 g_i=3m+(i\bmod m),qquad
 P_0=0,qquad P_j=\sum_{0\le i<j}g_i.
\]
The sum of one period and the right endpoint of the resulting chain are
\[
 S=\sum_{i=0}^{m-1}(3m+i)=\frac{m(7m-1)}2,qquad
 L=P_N=RS.
\]
By Bertrand's postulate, there is a prime~$p$ satisfying
\[
 4(D+L)+1<p<2\bigl(4(D+L)+1\bigr).
\]
Set
\[
 \mathcal P=\{D+P_0,\ldots,D+P_N\},qquad
 \mathcal Z=\mathcal P\cup(p-1-\mathcal P).
\]
The two displayed sets are disjoint because $p-1>4(D+L)$.  The definition
gives the reflection symmetry and excludes both endpoints.  All gaps within
either half are at least~$3m$, while the gap between the two halves is
$p-1-2(D+L)>2(D+L)$, so there are no consecutive points.  Moreover,
\[
 p>4(D+L)+1>\frac{m^2(7m-1)}{50}>144m^2=H^2.
\]

We next verify the pair cap.  Let
\[
 W(h)=\#\{(a,b):0\le a<b\le N,\ P_b-P_a=h\}.
\]
Every within-half difference is at most~$L$, and each occurs once in
each reflected half, whereas every cross-half difference is larger than
$p-1-2(D+L)>2(D+L)$.  Thus the two types have disjoint ranges.

For $1\le\ell<m$, define
\[
 F_\ell(t)=\sum_{j=0}^{\ell-1}((t+j)\bmod m),
 \qquad 0\le t<m.
\]
On $0\le t\le m-\ell$ this is a strictly increasing affine function
of~$t$, while on $m-\ell+1\le t<m$ it is a strictly decreasing affine
function.  Thus every value of $F_\ell$ has at most two preimages.
For a block of length~$\ell<m$,
\[
 P_{a+\ell}-P_a=3m\ell+F_\ell(a\bmod m).
\]
Each residue class modulo~$m$ contains at most~$R$ admissible starts,
so a fixed length~$\ell<m$ represents a fixed difference at most
$2R$ times.

Suppose that $1\le h<m^2$.  If $W(h)>0$, then the block length satisfies
\[
 \ell\le\frac{h}{3m}<\frac m3.
\]
Summing the preceding two-fiber bound over all possible lengths gives
\[
 W(h)\le2R\left\lfloor\frac{h}{3m}\right\rfloor.
\]
There are two reflected within-half copies, and hence, for $h\ge2$,
\[
 A_{\mathcal Z}(h)=2W(h)
 \le\frac{4Rh}{3m}=\frac h{75}\le3(h-1).
\]
For $h=1$ the count is zero.  If instead $m^2\le h\le L$, strict
monotonicity of the $P_a$ gives \(W(h)\le N\), whence
\[
 A_{\mathcal Z}(h)=2W(h)\le2N=\frac{m^2}{50}
 \le h\le3(h-1).
\]

It remains to consider a cross-half difference.  Such a difference has
the form
\[
 h=p-1-2D-P_a-P_b.
\]
For each~$a$, strict monotonicity gives at most one possible~$b$, so its
multiplicity is at most $N+1$.  Since $h>2(D+L)$, we have
\(h-1\ge2(D+L)\).  Moreover,
\(L=N(7m-1)/2\), and hence \(N+1\le2(D+L)\).  Therefore
\[
 A_{\mathcal Z}(h)\le N+1\le2(D+L)\le h-1\le3(h-1).
\]
All remaining differences have multiplicity zero.  This proves the pair
cap for every~$h$.

Finally, each half has $N+1$ points and therefore $N-2$ consecutive
four-point windows.  The span of any such window is the sum of three
successive gaps.  The largest possible sum is
\[
 (4m-3)+(4m-2)+(4m-1)=12m-6<H;
\]
wrapping across a period only makes it smaller.  No short window crosses
the central gap.  A pair in the left half has sum at most
$2(D+L)<p-1$,
while a pair in the right half has sum at least
$2(p-1-D-L)>p-1$.  Thus all $2(N-2)$ windows are off-center, and
\[
 2(N-2)=\frac{m^2}{50}-4=\frac{H^2}{7200}-4.
\]
Varying~$m$ through sufficiently large multiples of~$100$ completes the
construction.
\end{proof}

\begin{corollary}[Local set constraints do not control labeled separated energy]
\label{cor:pair-cap-separated-energy}
For every sufficiently large odd prime~$p$, there is a set
$\mathcal Z_p\subseteq\{0,\ldots,p-1\}$ satisfying properties
\textnormal{(1)--(3)} of Proposition~\ref{prop:pair-cap-extremal} such that,
if all points of $\mathcal Z_p$ are artificially assigned one color, then
the following hold.  Among the consecutive four-point windows of that color,
take those
whose three gaps are at least~$2$, whose gap triple is not constant, which
have no reflection-centered adjacent pair, and whose span is at most
$\lfloor\sqrt p\rfloor$.  After quotienting these windows by reflection,
their ordered separated energy is
\[
 E_p^{\mathrm{sep}}\gg p^{4/3}.
\]
Here two reflection orbits are separated when none of the two intervals in
the reflection-invariant support of one meets either interval in the support
of the other.  Consequently these combinatorial properties alone permit
\[
 \sum_{X<p\le2X}E_p^{\mathrm{sep}}
 \gg\frac{X^{7/3}}{\log X},
\]
which exceeds the weak projective-dispersion scale
\[
\frac{X^{2+o(1)}}{\log X}.
\]
This is a logical countermodel for arguments using only the listed
set-theoretic constraints: the artificial color is not asserted to be a
projective state of the Ap\'ery cocycle.  In particular, the construction
does not give a lower bound for the separated energy of the actual Ap\'ery
orbit.
\end{corollary}

\begin{proof}
For each sufficiently large~$p$, choose
\[
 m=100\left\lfloor\frac{p^{1/3}}{1000}\right\rfloor.
\]
Thus $m$ is a multiple of~$100$, $m\asymp p^{1/3}$, and $m\ge1100$.
Use the notation and construction in the proof of
Proposition~\ref{prop:pair-cap-extremal}, except that~$p$ is now prescribed:
\[
 D=1,\qquad R=\frac m{100},\qquad N=mR,\qquad
 L=\frac{m^2(7m-1)}{200}.
\]
Since $m\le p^{1/3}/10$,
\[
 4(D+L)<4+\frac{7m^3}{50}<p-1
\]
for all sufficiently large~$p$.  Hence the same two-packet construction
\[
 \mathcal Z_p
 =\{D+P_0,\ldots,D+P_N\}
  \cup\{p-1-D-P_0,\ldots,p-1-D-P_N\}
\]
is disjoint, and the pair-cap proof applies verbatim.  In particular,
$\mathcal Z_p$ is reflection invariant, omits both endpoints, has no
adjacent points, and satisfies
\[
 A_{\mathcal Z_p}(h)\le3(h-1)
 \qquad(1\le h<p).
\]

Each packet has
\[
 k=N-2=\frac{m^2}{100}-2
\]
consecutive four-point windows.  Reflection pairs the $k$ windows in the
left packet bijectively with the $k$ windows in the right packet, so there
are exactly $k$ reflection orbits.  Every gap is at least~$3m$.  Three
successive gaps are either three successive members of
\[
 3m,3m+1,\ldots,4m-1
\]
or one of the two triples crossing the end of this word; in every case the
three gaps are not all equal.  The span is at most $12m-6$, which is below
$\sqrt p$ for all sufficiently large~$p$.  The proof of the proposition
also shows that every such window is off-center.  The three global windows
that cross the gap between the packets have span greater than
$p-1-2(D+L)>2(D+L)$ and are therefore not selected.

Index the selected left-packet windows by $0,\ldots,k-1$.  Two such windows
have intersecting closed intervals exactly when their indices differ by at
most~$3$.  The two packets are separated by the central gap, so the same
criterion is equivalent to intersection of their full reflection-invariant
supports.  The ordered quotient overlap count is consequently
\[
 2\bigl((k-1)+(k-2)+(k-3)\bigr)=6k-12.
\]
There are $k(k-1)$ ordered distinct orbit pairs in total, and hence
\[
 E_p^{\mathrm{sep}}
 =k(k-1)-(6k-12)
 =(k-3)(k-4)
 \asymp m^4\asymp p^{4/3}.
\]
Applying this construction separately to every prime in $(X,2X]$ and using
the prime number theorem gives the displayed dyadic lower bound for the
artificially labeled model.  What fails for the genuine projective orbit is
precisely the unproved assertion that distinct packet windows carry the same
Ap\'ery projective state.
\end{proof}
